% =====================================================================
%  PRÉAMBULE COMMUN — Carnet d'entraînement, Fiche 2 (Nomenclature)
%  (inséré physiquement dans chaque document pour qu'il soit autonome)
% =====================================================================
\documentclass[11pt,a4paper]{article}

% ---------- Encodage & langue ----------
\usepackage[utf8]{inputenc}
\usepackage[T1]{fontenc}
\usepackage{lmodern}
\usepackage[french]{babel}

% ---------- Mathématiques & chimie ----------
\usepackage{amsmath,amssymb,mathtools}
\usepackage{icomma}
\usepackage[version=4]{mhchem}
\usepackage{siunitx}
\sisetup{output-decimal-marker={,},per-mode=symbol}

% ---------- Graphismes & mise en page ----------
\usepackage{xcolor}
\usepackage{colortbl}
\usepackage{tikz}
\usepackage[most]{tcolorbox}
\usepackage{enumitem}
\usepackage{array}
\usepackage{booktabs}
\usepackage{multicol}
\usepackage{fancyhdr}
\usepackage{titlesec}
\usepackage[hidelinks]{hyperref}
\usetikzlibrary{arrows.meta,positioning,calc,shapes.geometric,shapes.misc,
  fit,backgrounds,decorations.pathreplacing}
\usepackage[a4paper,margin=2.2cm,headheight=28pt,headsep=10pt]{geometry}

% =====================================================================
%  PALETTE (charte « Chimie » — mauve/prune du cours)
% =====================================================================
\definecolor{primary}{RGB}{146,95,161}
\definecolor{primarydark}{RGB}{92,55,108}
\definecolor{defcol}{RGB}{0,114,178}
\definecolor{thmcol}{RGB}{0,140,120}
\definecolor{formulacol}{RGB}{198,93,9}
\definecolor{remarkcol}{RGB}{120,94,200}
\definecolor{attncol}{RGB}{200,40,55}
\definecolor{excol}{RGB}{0,135,90}
\definecolor{methcol}{RGB}{90,90,100}
\definecolor{metalcol}{RGB}{198,150,40}
\definecolor{nonmetcol}{RGB}{0,140,150}
\definecolor{oxycol}{RGB}{210,55,45}
\definecolor{hydcol}{RGB}{60,120,210}
\definecolor{catcol}{RGB}{200,70,120}
\definecolor{ancol}{RGB}{40,110,190}
\definecolor{saltcol}{RGB}{120,90,160}

% =====================================================================
%  STYLE COMMUN DES BOÎTES
% =====================================================================
\tcbset{
  boxbase/.style={enhanced,breakable,boxrule=0.9pt,arc=2.5pt,
     left=3mm,right=3mm,top=2.3mm,bottom=2.3mm,
     fonttitle=\bfseries,coltitle=white,
     toptitle=0.6mm,bottomtitle=0.6mm},
}

\newtcolorbox{definition}[1][]{%
  boxbase,colback=defcol!5,colframe=defcol,
  title={\textsc{Définition}\ifx\relax#1\relax\relax\else\textmd{ --- \itshape #1}\fi}}
\newtcolorbox{regle}[1][]{%
  boxbase,colback=thmcol!6,colframe=thmcol,
  title={\textsc{Règle}\ifx\relax#1\relax\relax\else\textmd{ --- \itshape #1}\fi}}
\newtcolorbox{formule}[1][]{%
  boxbase,colback=formulacol!6,colframe=formulacol,
  title={\textsc{Méthode-clé}\ifx\relax#1\relax\relax\else\textmd{ --- \itshape #1}\fi}}
\newtcolorbox{remarque}[1][]{%
  boxbase,colback=remarkcol!6,colframe=remarkcol,
  title={\textsc{Remarque}\ifx\relax#1\relax\relax\else\textmd{ : \itshape #1}\fi}}
\newtcolorbox{attention}[1][]{%
  boxbase,colback=attncol!6,colframe=attncol,
  title={\textsc{Attention}\ifx\relax#1\relax\relax\else\textmd{ : \itshape #1}\fi}}
\newtcolorbox{exemple}[1][]{%
  boxbase,colback=excol!5,colframe=excol,
  title={\textsc{Exemple}\ifx\relax#1\relax\relax\else\textmd{ : \itshape #1}\fi}}
\newtcolorbox{methode}[1][]{%
  boxbase,colback=methcol!7,colframe=methcol,sharp corners,
  title={\textsc{Méthode}\ifx\relax#1\relax\relax\else\textmd{ : \itshape #1}\fi}}
\newtcolorbox{astuce}[1][]{%
  boxbase,colback=primary!6,colframe=primary,
  title={\textsc{Astuce concours}\ifx\relax#1\relax\relax\else\textmd{ : \itshape #1}\fi}}

% ---- Exercice (numéroté) ----
\newtcolorbox[auto counter]{exercice}[1][]{%
  boxbase,colback=primary!4,colframe=primarydark,
  title={\textsc{Exercice}~\thetcbcounter\ifx\relax#1\relax\relax\else\textmd{ --- \itshape #1}\fi}}

% ---- Corrigé (numéroté) ----
\newtcolorbox[auto counter]{corrige}[1][]{%
  boxbase,colback=excol!4,colframe=excol,
  title={\textsc{Corrigé}~\thetcbcounter\ifx\relax#1\relax\relax\else\textmd{ --- \itshape #1}\fi}}

% =====================================================================
%  EN-TÊTES ET PIEDS DE PAGE
% =====================================================================
\pagestyle{fancy}
\fancyhf{}
\fancyhead[L]{\small\textcolor{primarydark}{\textbf{Fiche 2} — Nomenclature}}
\fancyhead[R]{\small\textcolor{primarydark}{\textsc{Carnet d'entraînement}}}
\fancyfoot[C]{\small\thepage}
\renewcommand{\headrulewidth}{0.8pt}
\renewcommand{\headrule}{\hbox to\headwidth{\color{primary}\leaders\hrule height \headrulewidth\hfill}}

\titleformat{\section}{\Large\bfseries\color{primarydark}}{\thesection}{0.6em}{}
\titleformat{\subsection}{\large\bfseries\color{primary}}{\thesubsection}{0.6em}{}
\titleformat{\subsubsection}{\normalsize\bfseries\color{primary!80!black}}{}{0em}{}

% ---- Bandeau de titre réutilisable : \bandeau{Thème}{Sous-titre} ----
\newcommand{\bandeau}[2]{%
  \begin{center}
  \begin{tikzpicture}
  \node[rectangle,rounded corners=7pt,fill=primary,text=white,
        minimum width=16.0cm,minimum height=1.7cm,align=center,inner sep=8pt]
    {{\Large\bfseries #1}\\[3pt]{\normalsize #2}};
  \end{tikzpicture}
  \end{center}
  \vspace{0.3cm}}
\begin{document}
\bandeau{Thème 2 — Facile}{Du nom à la formule : croisement des valences --- Corrigé}

\begin{corrige}[croiser deux ions monoatomiques]
\begin{enumerate}[label=\alph*),leftmargin=1.8em,topsep=1pt,itemsep=1pt]
  \item \ce{Na+}$+$\ce{Cl-} $\to$ \ce{NaCl} (1 et 1).
  \item \ce{Ca^2+}$+$\ce{O^2-} $\to$ \ce{Ca2O2} $\Rightarrow$ simplifié \ce{CaO}.
  \item \ce{Al^3+}$+$\ce{O^2-} $\to$ \ce{Al2O3} (croisement 2 et 3, déjà simplifié).
  \item \ce{Mg^2+}$+$\ce{Cl-} $\to$ \ce{MgCl2}.
  \item \ce{Na+}$+$\ce{S^2-} $\to$ \ce{Na2S}.
\end{enumerate}
\end{corrige}

\begin{corrige}[où mettre les parenthèses ?]
\begin{enumerate}[label=\alph*),leftmargin=1.8em,topsep=1pt,itemsep=1pt]
  \item \ce{Ca(OH)2} (indice 2 $\geq$ 2 $\Rightarrow$ parenthèses).
  \item \ce{NaNO3} (indice 1 $\Rightarrow$ pas de parenthèses).
  \item \ce{Al(NO3)3} (indice 3 $\Rightarrow$ parenthèses).
  \item \ce{(NH4)2SO4} (le cation \ce{NH4} est polyatomique et présent 2 fois $\Rightarrow$ parenthèses ;
        \ce{SO4} est présent 1 fois, sans parenthèses).
  \item \ce{Ca3(PO4)2} (croisement 3 et 2 ; \ce{PO4} $\times 2 \Rightarrow$ parenthèses).
\end{enumerate}
\end{corrige}

\begin{corrige}[simplifier les indices]
\begin{enumerate}[label=\alph*),leftmargin=1.8em,topsep=1pt,itemsep=1pt]
  \item \ce{Ca2O2} $\Rightarrow$ \ce{CaO} (diviser par 2).
  \item \ce{Sn4O2}\,? Non : on croise \ce{Sn^4+} et \ce{O^2-} $\to$ \ce{Sn2O4} $\Rightarrow$ \ce{SnO2}
        (diviser par 2). \emph{(Le croisement place la valence 2 de l'oxygène sur l'étain et la valence 4
        de l'étain sur l'oxygène.)}
  \item \ce{Mg2S2} $\Rightarrow$ \ce{MgS} (diviser par 2).
\end{enumerate}
\end{corrige}

\begin{corrige}[du nom à l'ion et sa valence]
\begin{center}
\begin{tabular}{lcc}
\toprule
Nom & Formule de l'ion & Valence \\
\midrule
nitrate & \ce{NO3-} & 1 \\
sulfate & \ce{SO4^2-} & 2 \\
phosphate & \ce{PO4^3-} & 3 \\
hydroxyde & \ce{OH-} & 1 \\
carbonate & \ce{CO3^2-} & 2 \\
ammonium & \ce{NH4+} & 1 (cation) \\
\bottomrule
\end{tabular}
\end{center}
\end{corrige}

\begin{corrige}[l'ammonium joue le rôle d'un métal]
\begin{enumerate}[label=\alph*),leftmargin=1.8em,topsep=1pt,itemsep=1pt]
  \item \ce{NH4Cl} (1 et 1). \qquad b) \ce{NH4NO3} (1 et 1).
  \item[c)] \ce{(NH4)2CO3} (croisement 2 et 1 ; \ce{NH4} $\times 2 \Rightarrow$ parenthèses).
  \item[d)] \ce{(NH4)3PO4} (croisement 3 et 1 ; \ce{NH4} $\times 3 \Rightarrow$ parenthèses).
\end{enumerate}
\end{corrige}
\end{document}
